\section{Depth-Dependent Correlates of Recovery}
\label{sec:observation}

\subsection{Likelihood and capability are different observables}
Perplexity aggregates token-level negative log-likelihood over a corpus. A
knowledge-sensitive task instead asks whether a particular answer can be
selected or generated under a task interface. Healing can therefore improve the
former without changing the latter at the same rate. This is not a claim that
perplexity is uninformative; it is a claim that it is not a certificate for every
behavior after a structural intervention.

\subsection{Readability changes across depth}
% OLMo same-model probe saturation depths.
% Sources: results/probe_linguistic_olmo2_7b.json and
% results/knowledge_logit_lens_OLMo-2-1124-7B.json.
\begin{table}[t]
\centering
\scriptsize
\setlength{\tabcolsep}{4pt}
\begin{tabular}{@{}lcc@{}}
\toprule
\textbf{OLMo-2-7B quantity} & \textbf{layer} & \textbf{fractional depth} \\
\midrule
semantic probe mean sat$_{95}$ & 2.33 & $0.073L$ \\
MMLU onset (chance$+0.05$) & 18 & $0.562L$ \\
MMLU sat$_{95}$ & 19 & $0.594L$ \\
next-token sat$_{95}$ & 32 & $1.000L$ \\
\bottomrule
\end{tabular}
\caption{\textbf{Three ordered readout depths on the same OLMo model.}
Semantic sat$_{95}$ averages WiC, SST-2, and RTE depths measured relative to
each final-layer probe score. MMLU uses the final norm and output head;
next-token sat$_{95}$ is exact top-1 agreement with the final layer. These are
readout thresholds, not causal boundaries.}
\label{tab:mechanism}
\end{table}


On the intact OLMo-2-7B model, a final-head logit lens keeps MMLU answer letters
near chance through layer 17 and rises sharply at layers 18--19. Frozen semantic
classifiers become readable much earlier, while exact next-token agreement with
the final layer saturates only at the top. Figure~\ref{fig:twodepths} summarizes
this probe-dependent ordering.

\begin{figure}[t]
\centering
\includegraphics[width=\columnwidth]{figures/fig_two_depths.pdf}
\caption{\textbf{Probe-dependent readout depths on OLMo-2-7B.} The shaded band
marks prefix cuts used by the available behavioral frontier. Thresholds use
different readouts and are descriptive; they are neither directly comparable
causal boundaries nor locations of stored knowledge.}
\label{fig:twodepths}
\end{figure}

These probe-dependent readout depths are descriptive correlates, not causal
boundaries. They motivate measuring recovery at multiple behavioral levels
rather than assuming a single depth threshold, but they do not predict which
capabilities will return after pruning.

\subsection{Implications for interpreting recovery}
The observations clarify three choices in the preceding study design. First,
recovery is tracked along training because a single endpoint cannot distinguish
delayed from absent change. Second, answer-letter and complete-option scoring are
reported separately because they mix competence and readout differently. Third,
pruning policy is varied because preserving selected upper blocks can change the
result even at the same final depth. These are interpretation requirements, not
a causal theory of where capability resides.
